\section{The Proof Strategy}
\label{sec:proof-strategy}

Euclid's theorem is formalized as the following lemma:

\begin{equation*}
\forall\ \text{primes} \in \text{List[Prime]},\ \text{primes} \neq \emptyset \implies
\exists\ p \notin \text{primes} : \operatorname{isPrime}(p)
\end{equation*}

In the source, this is expressed by \texttt{PrimeProperties::euclidTheorem};
the verification reference is given in Section~3.4 and Appendix A.3.

\begin{itemize}
  \item Stage 1: $\operatorname{primorial}(L)+1$ is coprime to every prime
    in the list (Section~3.1).
  \item Stage 2: find a prime divisor of $\operatorname{primorial}(L)+1$
    via \texttt{findSmallestDivisor} (Section~3.2).
  \item Stage 3: the new prime is not in the original list (Section~3.3).
  \item Main theorem: combine stages 1--3 into Euclid's theorem
    (Section~3.4).
\end{itemize}

The proof proceeds in three stages:

\begin{enumerate}
  \item \textbf{Primorial-plus-one is not divisible by any prime in the
    list}: Show that $\operatorname{primorial}(L) + 1 \bmod p_i = 1 \neq 0$
    for every $p_i \in L$.
  \item \textbf{Smallest divisor is a new prime}: Find the smallest
    divisor $d > 1$ of $\operatorname{primorial}(L) + 1$. Prove $d$ is
    prime and is not in the original list.
  \item \textbf{Construct the result}: Return either $d$ (if
    $d < \operatorname{primorial}(L) + 1$) or
    $\operatorname{primorial}(L) + 1$ itself (if it is prime).
\end{enumerate}

\subsection{Stage 1: Primorial-plus-one Modulo All Primes}
\label{sec:stage-1}

The first stage is captured by the lemma \texttt{primorialPlusOneModAny}.
Let $L = [p_1,\dots,p_k]$ be the finite list of known primes and
$P=\operatorname{primorial}(L)$. For each $p_i \in L$, the product $P$
contains $p_i$ as one factor, so $P$ is divisible by $p_i$. Adding one
moves the residue from $0$ to $1$, and because every prime is greater than
$1$, that residue is nonzero.

\begin{equation*}
\begin{aligned}
\operatorname{primorial}(L)
  &= p_i \cdot \prod_{j\ne i} p_j
  &&\text{[By Definition]} \\
\operatorname{mod}(\operatorname{primorial}(L), p_i)
  &= 0
  &&\text{[Product Contains }p_i\text{]} \\
\operatorname{mod}(\operatorname{primorial}(L)+1, p_i)
  &= \operatorname{mod}(1, p_i)
  &&\text{[Modulo Shift]} \\
  &= 1
  &&\text{[Since }1 < p_i\text{]} \\
  &\ne 0
  &&\blacksquare\ \text{[Q.E.D.]}
\end{aligned}
\end{equation*}

The verified source proves this by induction over the list. At each step,
the current prime is split out of the primorial product, the divisibility
of the remaining product is preserved by multiplication, and the induction
hypothesis continues over the tail. The loop step is built from three
verified arithmetic properties.

\textbf{Small Dividend Remainder.} A nonnegative dividend smaller than the
divisor is already its own remainder. In the Euclid step, this gives both
$\operatorname{mod}(0,p)=0$ and $\operatorname{mod}(1,p)=1$ because $p>1$.

\begin{equation*}
\begin{aligned}
0 \le a < b
&\Rightarrow \operatorname{mod}(a,b)=a
&&\text{[Small Dividend]} \\
p>1
&\Rightarrow \operatorname{mod}(0,p)=0
\land \operatorname{mod}(1,p)=1
&&\text{[Substitution]}
\end{aligned}
\end{equation*}

This property is verified in
\href{https://github.com/thiagomata/prime-numbers/blob/master/src/main/scala/v1/chapter2/div/properties/ModSmallDividend.scala}{\texttt{ModSmallDividend::modSmallDividend}}.

\textbf{Zero Remainder Preserved by Multiplication.} If a number is
divisible by $b$, multiplying it by any nonnegative factor preserves
divisibility by $b$. This is the step that turns the explicit factor $p$
in the primorial into $\operatorname{mod}(p\cdot k,p)=0$.

\begin{equation*}
\begin{aligned}
\operatorname{mod}(a,b)=0
&\Rightarrow \operatorname{mod}(a\cdot m,b)=0
&&\text{[Multiplication Preserves Zero Remainder]} \\
\operatorname{mod}(p,p)=0
&\Rightarrow \operatorname{mod}(p\cdot k,p)=0
&&\text{[Substitution]}
\end{aligned}
\end{equation*}

This property is verified in
\href{https://github.com/thiagomata/prime-numbers/blob/master/src/main/scala/v1/chapter2/div/properties/AdditionAndMultiplication.scala}{\texttt{AdditionAndMultiplication::ATimesBSameMod}}.

\textbf{Adding One After a Multiple.} Once the primorial part is known to
be divisible by $p$, adding one gives the same remainder as one itself.
Together with the small-dividend property, this proves the Euclid number
has nonzero remainder modulo every original prime.

\begin{equation*}
\begin{aligned}
\operatorname{mod}(m,b)=0
&\Rightarrow \operatorname{mod}(m+c,b)=\operatorname{mod}(c,b)
&&\text{[Modulo Shift]} \\
\operatorname{mod}(p\cdot k,p)=0
&\Rightarrow \operatorname{mod}(p\cdot k+1,p)=\operatorname{mod}(1,p)=1
&&\text{[Substitution]} \\
&\Rightarrow \operatorname{mod}(p\cdot k+1,p)\ne0
&&\blacksquare\ \text{[Q.E.D.]}
\end{aligned}
\end{equation*}

This property is verified in
\href{https://github.com/thiagomata/prime-numbers/blob/master/src/main/scala/v1/chapter2/div/properties/ModOperations.scala}{\texttt{ModOperations::modZeroPlusC}}.

This property is verified in
\href{https://github.com/thiagomata/prime-numbers/blob/master/src/main/scala/v1/chapter5/prime/properties/PrimeProperties.scala}{\texttt{PrimeProperties::primorialPlusOneModAny}}.
A short public wrapper excerpt is included in Appendix A.1.

\subsection{Stage 2: Finding a New Prime}
\label{sec:stage-2}

Once we know $\operatorname{primorial}(L)+1$ is not divisible by any prime
in $L$, let

\begin{equation*}
\begin{aligned}
N &= \operatorname{primorial}(L)+1, \\
d &= \operatorname{findSmallestDivisor}(N,2).
\end{aligned}
\end{equation*}

There are two cases. If $d=N$, the divisor search found no proper divisor
in $[2,N)$, so $N$ is prime. If $d < N$, then $d$ divides $N$ and no
smaller integer greater than $1$ divides $N$. If $d$ were composite, it
would have a non-trivial divisor $e$ with $1 < e < d$; since $e$ divides
$d$ and $d$ divides $N$, $e$ would divide $N$, contradicting the
minimality of $d$. Hence $d$ is prime.

\begin{equation*}
\begin{aligned}
d=N
&\Rightarrow \forall e\in[2,N),\operatorname{mod}(N,e)\ne0.
&&\text{[No Proper Divisor Found]} \\
&\Rightarrow \operatorname{isPrime}(N)
&&\text{[Prime Definition]} \\
&d < N \land \operatorname{mod}(N,d)=0 \Rightarrow \operatorname{isPrime}(d)
&&\text{[Minimal Divisor]}
\end{aligned}
\end{equation*}

The construction of the new prime is verified in
\href{https://github.com/thiagomata/prime-numbers/blob/master/src/main/scala/v1/chapter5/prime/properties/PrimeProperties.scala}{\texttt{PrimeProperties::newPrimeFromEuclid}}.
A short public wrapper excerpt is included in Appendix A.2.

\subsection{Stage 3: Proving the New Prime is Not in the List}
\label{sec:stage-3}

The final and most subtle step is proving that the newly found prime $d$
(or $N$ itself) is \textbf{not} in the original list.

Let $v$ be the divisor chosen in Stage 2: either $v=N$ when $N$ is prime,
or $v=d$ when $d$ is the smallest proper divisor of $N$. In both cases,
$\operatorname{mod}(N,v)=0$. Now take any prime $p$ from the original
list. Because $N=\operatorname{primorial}(L)+1$, the same argument from
Stage 1 gives $\operatorname{mod}(N,p)=1$. If $p=v$, then $N$ would have
two incompatible remainders modulo the same positive divisor: $0$ and $1$.
Therefore no element of $L$ equals $v$.

\begin{equation*}
\begin{aligned}
N &= \operatorname{primorial}(L) + 1
&&\text{[Euclid Construction]} \\
  &= p \cdot k + 1
&&\text{[Unfold Product at }p\text{]} \\
\operatorname{mod}(N,p)
  &= \operatorname{mod}(p\cdot k+1,p) \\
  &= \operatorname{mod}(1,p)
&&\text{[Multiple of }p\text{ Drops Out]} \\
  &= 1
&&\text{[Since }p>1\text{]} \\
\operatorname{mod}(N,v)
  &= 0
&&\text{[Chosen Divisor]} \\
p=v
  &\Rightarrow 1=0
&&\text{[Contradiction]} \\
\therefore\ p &\ne v
&&\blacksquare\ \text{[Q.E.D.]}
\end{aligned}
\end{equation*}

This non-membership argument is verified by the private helper
\texttt{euclidTailLoop}, which establishes
\texttt{valueNotMatchesAny(primes, v)} for the chosen divisor $v$ in
\href{https://github.com/thiagomata/prime-numbers/blob/master/src/main/scala/v1/chapter5/prime/properties/PrimeProperties.scala}{\texttt{PrimeProperties}}.

\subsection{The Main Theorem}
\label{sec:main-theorem}

The main theorem combines the primorial-plus-one lemma, smallest-divisor
primality, and non-membership argument. If $N$ is prime, then $N$ itself
is the new prime. Otherwise, the smallest divisor $d$ of $N$ is prime and
cannot belong to the original list.

\begin{equation*}
\begin{aligned}
L &\ne [] \\
N &= \operatorname{primorial}(L)+1 \\
d &= \operatorname{findSmallestDivisor}(N,2) \\
d=N
&\Rightarrow \operatorname{isPrime}(N)\land N\notin L
&&\text{[Stages 2 and 3]} \\
d < N
&\Rightarrow \operatorname{isPrime}(d)\land d\notin L
&&\text{[Stages 2 and 3]}
\end{aligned}
\end{equation*}

\begin{equation*}
\therefore\ \exists p:\operatorname{isPrime}(p)\land p\notin L
\quad \text{[Case Split]}\quad\blacksquare\ \text{[Q.E.D.]}
\end{equation*}

This property is verified in
\href{https://github.com/thiagomata/prime-numbers/blob/master/src/main/scala/v1/chapter5/prime/properties/PrimeProperties.scala}{\texttt{PrimeProperties::euclidTheorem}}.
The public theorem wrapper is shown in Appendix A.3.
